\chapter{Notation and Symbols}\label{app:notation}

This appendix collects the principal symbols used in the manual. The sign and unit
conventions they obey are stated in the front-matter section ``Notation and conventions''.

\begin{longtable}{@{}l p{0.72\linewidth}@{}}
\toprule
Symbol & Meaning \\
\midrule \endhead
$\pos=(\lonx,\latx,\depx)$ & Particle position: longitude, latitude (degrees), depth (m, positive down). \\
$\velvec=(u,v,\Wvel)$ & Velocity: horizontal components (m\,s$^{-1}$) and vertical velocity $\Wvel$ (m\,s$^{-1}$, positive down). \\
$t,\ \dt$ & Time and the (positive) integration step. \\
$\tdir$ & Integration-direction factor, $+1$ forward or $-1$ backward (the key \conf{time\_dir}). \\
$\Mlat$ & Metres per degree of latitude, $\SI{111120}{\metre\per\degree}$. \\
$\zzero$ & Metres per degree of longitude, $\max(\cos\latx\,\Mlat,1)$. \\
$\bwt{i}$ & Barycentric / Wachspress interpolation weight at vertex $i$. \\
$\interpfrac$ & Interpolation fraction, in time (\cref{eq:temp-linear}) or in depth (\cref{eq:vert-alpha}). \\
$N,\ E,\ K$ & Number of mesh nodes, elements, and vertical interfaces ($K-1$ layers). \\
$z_k$ & Depth of the $k$-th vertical interface. \\
$\sigma_k$ & Sigma value of the $k$-th interface (FESOM-C); $\sigma=1$ surface, $\sigma=0$ floor. \\
$H$ & Sigma-mesh water-column thickness, $\max(d+\eta,D_{\min})$. \\
$d,\ \eta$ & Still-water depth and sea-surface height. \\
$D_{\min},\ z_{0b}$ & Minimum column thickness and bottom clamp offset (sigma; keys \conf{Dmin}, \conf{z0b\_min}). \\
$K_h$ & Horizontal eddy diffusivity (m$^2$\,s$^{-1}$); $\sigma^2=2K_h t$. \\
$R$ & Earth radius of the mesh (\conf{mesh.r\_earth}), used in the local tangent-plane projection. \\
$R_M$ & Metric radius of the displacement, $\Mlat\cdot180/\pi$, used by the great-circle kicks (\cref{sec:diff-sphere}). \\
$\xi,\ \xi_1,\ \xi_2$ & Standard normal variates of the stochastic kicks (\cref{ch:diffusion}). \\
$A_z$ & Vertical eddy diffusivity (m$^2$\,s$^{-1}$), \cref{sec:diff-vertical}. \\
$c_k$ & Fraction of the step at which Runge--Kutta stage $k$ samples the flow. \\
$p_k,\ M$ & Data points of a depth-linear horizontal velocity (layer centres or levels; on a sigma mesh layer centres in $s$) and their number (\cref{sec:vert-uvdepth}). \\
$\ell_c,\ d_c,\ D_c$ & Last usable data point, bed depth and reader-declared data prefix of column $c$ (a node or an element), \cref{eq:vert-uv-last}. \\
$s$ & Sigma fraction $z/H_p$ of the particle's column, clipped to $[0,1]$ (\cref{eq:vert-uv-sigma}). \\
$\bwt{i}^{+}$ & Barycentric weight counted as zero when negative, in the wall rule (\cref{eq:vert-uv-bary}). \\
$\alpha_0,\ \alpha_k$ & Bracket fraction of a lazy blend and stage fraction of a stage blend (\cref{ch:temporal}). \\
$A_k,\ C_k$ & Doubled fan- and corner-triangle areas in the Wachspress weights (\cref{eq:wachspress}). \\
\bottomrule
\end{longtable}
